\lista{
  \textbf{Interfaz de usuario:} Gestiona barras de herramientas, progreso de carga, preferencias y se comunica con el sistema operativo. |
  \textbf{Motor del buscador:} Carga urls, soporta navegación básica (anterior/siguiente) y gestiona alertas de javascript. |
  \textbf{Motor de renderizado:} Presenta visualmente el recurso, interpretando html, xml, css y contenido embebido. |
  \textbf{Subsistema de comunicaciones:} Implementa protocolos de transferencia de archivos (ej. http). |
  \textbf{Intérprete de javascript:} Analiza y ejecuta código javascript, incluyendo eventos de ratón o teclado. |
  \textbf{Parser xml:} Carga la representación en árbol (dom) del documento html en memoria para facilitar el acceso a sus elementos. |
  \textbf{Componente de visualización:} Ofrece funcionalidades específicas para la presentación del contenido html. |
  \textbf{Subsistema de persistencia de datos:} Almacena historial, sesiones de usuario, preferencias, certificados de seguridad y cookies.
}
\nota{Para el examen}{
Diferencie entre el motor del buscador y el motor de renderizado. El primero gestiona la carga de la dirección, mientras que el segundo se encarga de la presentación visual del contenido interpretado.
}
